\documentclass[fontsize=11.5pt]{scrartcl}

\usepackage[letterpaper,hmargin=1.3in,vmargin=0.8in]{geometry}
\usepackage{microtype}
\usepackage{enumitem}
\usepackage[draft]{hyperref}
\usepackage{parskip}
\usepackage{graphicx}

\setlength{\parskip}{4pt plus 1pt}

\begin{document}
\pagestyle{empty}

\begingroup
\small
\setlength{\leftskip}{1.5em}% one paragraph-indent width; adjust as desired
Title:
Living Alignment

Authors:\par
{\setlength{\parskip}{0pt}\everypar{\hangindent=1.5em\hangafter=1}% indent wrapped affiliation lines
Zeb Kurth-Nelson, Prefrontal, London, UK\par
Ruairidh Battleday, Prefrontal, London, UK\par
James Whittington, Prefrontal, London, UK\par
Nenad Tomašev, Google DeepMind, London, UK\par
Erman Misirlisoy, Independent\par
Joel Z. Leibo, Google DeepMind, London, UK\par
Matthew M. Nour, Microsoft AI, London, UK; Department of Psychiatry, University of Oxford, Oxford, UK; Max Planck UCL Centre for Computational Psychiatry and Ageing, University College London, London, UK\par
Marc Guitart-Masip, Karolinska Institutet, Stockholm, Sweden\par
Steve Sullivan, Oregon Health and Science University, Portland, OR\par}

% Keywords: AI safety, open-endedness, evolutionary biology, ecosystems

% Keywords: AI alignment, AI safety, inverse RL, multi-agent, reward hacking, diversity, complexity, open-endedness, niche creation, evolution, genetic recombination, functional fixedness, cognitive control, pattern separation, ecosystems, regulatory frameworks, sociotechnical alignment, polycentric governance, inequality

\endgroup

\vspace{1em}

Dear PNAS Editors,

We submit our manuscript, \textit{Living Alignment}, for consideration as a Perspective in PNAS, following the Editorial Board's invitation of August 10, 2026 in response to our presubmission inquiry (tracking number 2026-22180PPER). Our paper offers a new perspective on the problem of AI alignment, inspired by life.

In the first half of the paper, we identify a common principle that has been discovered separately in multiple disciplines across the life and social sciences: \textbf{living systems are healthy and flourishing when they avoid fixation on particular modes of interaction}. By avoiding fixation, life maintains the potential to continually invent new solutions and patterns of organization. We walk through examples from a range of domains, including evolutionary biology, cell biology, neuroscience, cognitive science, social science and ecology. 

The second half of the paper applies the same principle to AI alignment. There is broad agreement that aligned AI systems should support futures that are healthy and flourishing. Therefore, we reason that \textbf{alignment means avoiding fixation on particular modes of interaction}. We use the problem of reward hacking to illustrate the message in detail, then discuss more generally what negative and positive outcomes look like in the living alignment framework.

If AI drives large changes in the world, this level of abstraction---where the science of life teaches us what it means for any kind of system to be healthy and flourishing---may generalize better to new patterns of organization. 

We are excited about submitting this article to PNAS. The article integrates theory and data from multiple fields and draws heavily on work published in PNAS. Our specialities span neuroscience, cognitive science, AI research, social science, physiology and psychiatry. The simple premise of our article is important and previously unarticulated, and it is sure to engage debate.

\vspace{1em}


\begin{minipage}{\textwidth}
Sincerely,

\includegraphics[height=4em]{../assets/signature.png}

Zeb Kurth-Nelson\\
on behalf of the authors
\end{minipage}


\end{document}
